\section{知的財産を分類するための二つの問い}
\label{theory:ip:scope}
知的財産の問題では、まず保護される対象を問い、次にどの制度で保護するかを問う。技術的な解決方法、文章や画像の表現、商品の出所を示す標識、秘密の製造情報は、同じ情報でも保護の根拠が異なる。情報は多人数が同時に利用できるため、物の占有と同じ方法で利用を制御できない。制度は創作・研究開発への誘因と成果の利用を調整する。

特許庁が産業財産権と総称するのは、特許権、実用新案権、意匠権、商標権の4種類である\parencite{ip_industrial_four}。著作権、回路配置利用権、育成者権は知的財産に関係する権利であっても、この4種類に含まれない。営業秘密、商号、商品等表示の保護も広い知的財産の体系に位置付く\parencite{ip_basic_act,ip_broad_scope}。広い体系へ含めることは、全てが登録により絶対的な独占権を生むという意味ではない。

\subsection{集合の補集合と選択肢番号}
\label{theory:ip:count}
列挙された対象の集合を \(L\)、4種類の集合を \(I\) とすると、産業財産権に含まれない対象は \(L\setminus I\) である。平成30年問6は広い知的財産の範囲外の数、令和2年問5は産業財産権の範囲外の数を問う。同じ名称が並んでも、引く集合が違えば数が変わる。対象数を求めた後で、その数を示す選択肢番号に対応させる。

\section{特許：定義と権利取得の要件を分ける}
\label{theory:ip:patent}
自然法則を利用した技術的思想の創作のうち高度なものが、特許法上の発明である。発明に該当しても、直ちに特許を得るわけではない。産業上利用できること、出願前の公知・公用・刊行物等に照らした新規性、進歩性などを別に審査する。公の秩序、善良の風俗又は公衆の衛生を害するおそれがある発明は保護対象から除かれる\parencite{ip_patent_statute}。

\subsection{出願前公開と登録後公開の区別}
\label{theory:ip:novelty}
制度が技術の公開と利用を促すことと、出願前に既に公表されていることを取得要件にすることは別である。通常、出願前公開は新規性に不利に働く。ただし公開された発明が例外なく永久に特許を得られない、とも断定しない。新規性喪失の例外には対象行為、期限、手続の条件がある。

平成30年改正では例外期間が6か月から1年へ延長され、2018年6月9日以後の出願への適用に公開日の経過措置も置かれた\parencite{ip_novelty_grace2018}。令和元年再試験の判断は、この当時すでに例外が存在したことと整合させる。例外は「刊行物で発表済みでなければならない」という要件を作らない。

\section{著作権：表現、創作性、権利の発生}
\label{theory:ip:copyright}
著作物は思想又は感情を創作的に表現したものである。学術的な事実や抽象的なアイデアそのものと、その説明文・図の具体的な表現を区別する。創作性は芸術的な傑作という評価を要求するものではないが、法の定義から削除されたわけでもない。

著作者人格権と著作権の享有に申請や審査を要求しない仕組みを無方式主義という\parencite{ip_copyright_statute}。特許庁へ登録しなければ著作権が発生しないという説明は、産業財産権の手続との混同である。権利の発生と、一定事項の登録制度が存在することも別である。

\subsection{人格的利益と財産的利用}
\label{theory:ip:personality}
公表権、氏名表示権、同一性保持権は著作者の人格的利益に関わる。複製、公衆送信などの利用を管理する財産的な権利と区別する。歌手が歌唱する立場にあることだけで、他者が作った歌詞を自由に変更する権限は生じない。利用許諾があることも、全ての改変への同意があることを当然には意味しない。

\subsection{個人の趣味と公開送信}
\label{theory:ip:public_web}
個人的な趣味であっても、公開Webページへの画像掲載は私的な複製と同じ範囲ではない。写真集の画像について、許諾や引用等の適法な根拠を欠く掲載を、趣味という目的だけで正当化しない。写真の権利と被写体に関わる別の権利も分けて考える。

\section{引用：出典表示だけでは条件を満たさない}
\label{theory:ip:quotation}
32条1項の引用は、公表された著作物を、公正な慣行と引用目的上正当な範囲に従って利用する仕組みである\parencite{ip_copyright_statute}。自分の論述に必要な部分を区別し、利用量と目的の関係を説明できることが重要になる。48条の出所明示も確認するが、出典を書けばどの量でも転載できるという制度ではない。

権利制限を満たす引用は、保護期間内であっても許諾なしにできる。保護期間満了後の利用、権利者の許諾に基づく利用、法が認める引用は異なる根拠である。「無断なら例外なく違法」という断定を避ける。

\subsection{官公資料の範囲と禁止表示}
\label{theory:ip:government}
法令など13条に掲げる対象が著作権の対象から除かれることと、官公庁が作成した全資料を自由に使えることは別である。32条2項は、国等が一般への周知を目的に自らの名義で公表した広報・調査統計・報告等を、説明材料として新聞紙、雑誌その他の刊行物へ転載する仕組みを定めるが、転載禁止の表示がある場合を除く。対象、目的、媒体、禁止表示を確認する\parencite{ip_copyright_statute}。

\subsection{引用のための翻訳と自由な翻案}
\label{theory:ip:translation}
適法な引用に対応する範囲では翻訳利用が認められ、外国語の原文をそのまま使うことだけが唯一の方法ではない。これは他者の文章を自由に翻案・改変してよいという意味ではない。取得した現在の条文では47条の6が翻訳等を扱う。平成25・26年度の問題について当時の条番号を現在番号に置き換えた証跡とはしない。二問の当時版全文の照合は未完として残す。

\section{営業秘密：情報の種類と法的成立を分ける}
\label{theory:ip:secret}
不正競争防止法2条6項の営業秘密は、秘密管理性、有用性、非公知性を同時に満たす情報である\parencite{ip_unfair_statute}。顧客名簿や新規事業計画はその典型になり得るが、資料名や社内で作成された事情だけで営業秘密になるわけではない。秘密として管理され、事業に有用で、公然と知られていないことが必要である。

\subsection{有用性と現実の使用}
\label{theory:ip:usefulness}
情報の有用性は、現に製品へ使われていることだけでは決まらない。将来の事業計画、失敗した実験の記録も、不要な研究費を避けるなどの価値を持ち得る。一方、脱税や有害物質の垂れ流し等の反社会的な活動の情報は、法が保護する正当な事業活動の有用性とは扱われない\parencite{ip_secret_handbook2017,ip_secret_commentary2018}。違法行為の情報を隠すことと、適法な技術情報を守ることを分ける。

\subsection{公開、秘密管理、漏洩経路}
\label{theory:ip:secrecy}
刊行物や公開された特許に記載された内容は、そのまま非公知の秘密とはできない。そこに記載されていない別のノウハウまで一緒に公開されたとは限らない。管理表示やアクセス制御は、対象者が秘密と認識できる管理の具体化である。漏洩経路には従業員、退職者、取引先、共同研究先、サイバー攻撃がある。経路の存在の確認と、特定年に急増したという量的主張の確認は別である。

\subsection{令和2年と令和3年の同じ省略表現}
\label{theory:ip:secret_ambiguity}
両年の設問は、顧客名簿等について、企業活動の過程で生まれたので営業秘密であると述べる。公式表は典型的な該当例として真とする組合せである。一方、三要件を省略しない一般命題として読むと、作成された由来だけでは成立を保証できない。令和3年の独立初回判断はこの読みから4、公式正答は5であった。

本教材は公式値を保持するとともに、初回判断と省略条件を残す。公式と一致するよう法の要件を変更しない。令和2年も同じ条件省略があるため、公式値と一致した事実だけで曖昧さを解消済みにしない。

\section{非享受目的：利用の目的から権利制限を判断する}
\label{theory:ip:nonenjoyment}
30条の4は、表現された思想又は感情を自ら享受し又は他人に享受させることを目的としない場合に、必要な限度で利用を認める。ただし著作権者の利益を不当に害する場合を除く。技術試験、情報解析、人の知覚によらない計算処理は例示であり、「AIだから許される」という機器名の判断ではない\parencite{ip_flexible2019,ip_copyright_statute}。

\subsection{六つの例とそれぞれの境界}
\label{theory:ip:six_cases}
2019年文化庁資料は、AI学習用の収集・第三者提供、プログラムの解析、複製機器や和紙の性能試験、送り仮名等の調査、場所の3DCG化、全文検索用の複製を具体例に挙げる\parencite{ip_flexible2019}。それぞれ表現の鑑賞ではなく、解析・性能・事実情報等を目的とする例である。

AI学習では、既存作品の表現を享受させる目的の併存や、解析用に販売されるデータベースの市場への影響を別に確認する。性能試験の複製物を鑑賞用に頒布する行為まで認めることにはならない。写真から場所を再現する場合も、元写真の創作的な表現を3DCGで承継して享受させる場合は区別する。

検索のための内部複製と検索結果の表示も別である。書名・著者名等の表示を目的とする例と、本文の一部を表示する利用についての47条の5の条件を混同しない。プログラム解析が著作権法上許される場合でも、契約など別の問題が全て解決するわけではない。

\section{法の時点を保った学習の手順}
\label{theory:ip:historical}
4種類への分類、引用の条件、三要件の論理は共通に参照し、保護対象の拡大や条番号、例外の改正は出題時点を確認する。非享受目的の2019年施行の規定を、それ以前の問題にそのまま遡及しない。現在の条文の説明、当時の行政資料で確認した範囲、未取得の過去版全文を区別する。

正誤表では各文の判断を先に定め、原典の列順で5候補と照合する。否定、数、組合せの違いを保ち、同じ公式番号が続くことを同じ答えの内容とみなさない。
